\documentclass[11pt,a4paper]{article}
\usepackage[utf8]{inputenc}
\usepackage[spanish]{babel}
\usepackage{amsmath}
\usepackage{amsfonts}
\usepackage{geometry}
\usepackage{xcolor}
\usepackage{booktabs}
\usepackage{titlesec}

% Definición de márgenes y geometría
\geometry{
    a4paper,
    left=25mm,
    right=25mm,
    top=25mm,
    bottom=25mm
}

% Colores elegantes para los títulos
\definecolor{primary}{RGB}{31, 78, 121}    % Azul marino
\definecolor{secondary}{RGB}{90, 90, 90}   % Gris oscuro

% Configuración de títulos estilo reporte limpio
\titleformat{\section}
  {\color{primary}\normalfont\Large\bfseries}
  {\thesection}{1em}{}[{\color{primary}\titlerule[1pt]}]

\titleformat{\subsection}
  {\color{secondary}\normalfont\large\bfseries}
  {\thesubsection}{1em}{}

\setlength{\parindent}{0pt}
\setlength{\parskip}{8pt}

\begin{document}

% Encabezado del documento
\begin{center}
    {\color{primary}\LARGE\bfseries Formulario de prestaciones} \\[0.2cm]
    {\color{secondary}\large Daniel Rodulfo 31725032} \\[0.4cm]
    \rule{\linewidth}{0.5mm}
\end{center}

\vspace{0.5cm}

\section{Fórmulas Principales}


\subsection{1. Prestaciones (Rendimiento)}
\begin{equation}
    \text{Prestaciones} = \frac{1}{\text{Tiempo de Ejecución}}
\end{equation}

\subsection{2. Tiempo de Ejecución (Ecuación Básica)}
\begin{equation}
    \text{Tiempo de Ejecución} = \frac{\text{Ciclos de Reloj}}{\text{Frecuencia de Reloj}}
\end{equation}

\subsection{3. Ciclos de Reloj de la CPU}
\begin{equation}
    \text{Ciclos de Reloj} = \text{Número de Instrucciones} \times \text{CPI}
\end{equation}

\subsection{4. Ecuación reformulada del Rendimiento de CPU}

\begin{equation}
    \text{Tiempo de Ejecución} = \frac{\text{Número de Instrucciones} \times \text{CPI}}{\text{Frecuencia de Reloj}}
\end{equation}

\vspace{0.5cm}
\end{document}